\documentclass[10pt,landscape]{article}
\usepackage[margin=1.6cm]{geometry}
\usepackage[utf8]{inputenc}
\usepackage[T1]{fontenc}
\usepackage[spanish,es-noshorthands]{babel}
\usepackage{booktabs}
\usepackage{tabularx}
\usepackage{longtable}
\usepackage{array}
\usepackage{xcolor}
\usepackage{colortbl}
\usepackage{ragged2e}
\definecolor{prim}{HTML}{1B2A4A}
\definecolor{primlight}{HTML}{D6E4F0}
\definecolor{gris}{HTML}{F7F7F5}
\newcolumntype{L}[1]{>{\RaggedRight\arraybackslash}p{#1}}
\renewcommand{\arraystretch}{1.12}
\setlength{\aboverulesep}{0pt}
\setlength{\belowrulesep}{0pt}
\title{\textcolor{prim}{Planificación LAB019--021: réplica del Lote 2 (LAB004--006)}\\[2pt]\large Vendimia 2026 --- mosto natural, escala laboratorio}
\author{}
\date{20/09/2026 (v2: Carbodoseur + tiempos de procesamiento)}
\begin{document}
\maketitle
\thispagestyle{empty}
\vspace{-1.2cm}

\section{Objeto y contexto}
Se replica el lote LAB004--LAB006 (mosto natural \emph{Sauvignon blanc}, levadura Zymaflore X5, reactores R-2L de 2 L) del 6--15 de abril de 2026, cuya implementación de medición de CO$_2$ presentó fugas y fue descartada del análisis (\texttt{EXCLUDED\_BATCHES}). Con el sistema de CO$_2$ correctamente conectado y verificado, se repiten los perfiles de temperatura y nutrición originales para obtener una serie continua confiable sobre el mismo diseño experimental, y se agrega la medición de CO$_2$ disuelto con \textbf{Carbodoseur} en eventos clave (observabilidad del pool de CO2 disuelto del modelo).

\textbf{Correspondencia:} LAB019 = R-2L-01 (F1), Protocolo A, réplica de LAB004; LAB020 = R-2L-02 (F2), Protocolo C, réplica de LAB006; LAB021 = R-2L-03 (F3), Protocolo A, réplica de LAB009 (grupo del 17/04, sin registro en Maestro). Se mantienen 2 réplicas del protocolo A (F1/F3) y la C como fermentación de contraste (arranque frío). Nota: los tres experimentos originales tuvieron CO$_2$ descartado por QC, por eso son los que más ganan con esta réplica.

\section{Análisis de los experimentos originales}
De las hojas \texttt{Eventos\_operacion}, \texttt{Perfiles\_temperatura} y \texttt{LAB004..LAB006} del workbook homologado:

\begin{table}[h]\centering
\begin{tabular}{L{6.2cm} L{7.5cm} L{7.5cm}}
\toprule
\rowcolor{primlight}\textbf{Hito} & \textbf{Protocolo A (LAB004 / LAB009)} & \textbf{Protocolo C (LAB006)} \\
\midrule
Pulso nutricional 2 (densidad 1040) & t=51 h (LAB004) / t=48 h vía ICS (LAB009, sin registro) & t=75 h (jue 09/04 16:00) \\
Cambio de T 1 & 21\,\textdegree C a t=99 h (LAB004) / t=96 h (LAB009) & 18\,\textdegree C a t=123 h \\
Cambio de T 2 & 16\,\textdegree C a t=171 h (LAB004) / t=168 h (LAB009, inferido) & 21\,\textdegree C a t=219 h \\
Fermentación seca (GF$<$4 g/L) & t$\approx$122--146 h & t$\approx$170--194 h \\
Última muestra del original & t=163 h (LAB004) / t=160 h (LAB009) & t=211 h (15 muestras) \\
Duración total & $\approx$163 h ($\approx$6.8 días) & $\approx$211 h ($\approx$8.8 días) \\
Grilla original & 2 muestras/día hábil (08:00 y 15:00) + sábado y domingo & ídem, extendida un día más \\
\bottomrule
\end{tabular}
\caption{Hitos de los experimentos originales (LAB004 y LAB006: siembra 06/04/2026 13:00; LAB009: siembra 17/04/2026 16:00).}
\end{table}

Con $t_0$ = lunes 21/09 a las 16:00: \textbf{A} seca entre el sábado 26 16:00 y el domingo 27 18:00 (t=120--146 h; LAB004 t=122--146, LAB009 t=120--144), fin del protocolo el lunes 28/09 en el día (t=160--171 h); \textbf{C} seca entre el lunes 28 y el martes 29/09 por la tarde (t=170--194 h), fin del protocolo el miércoles 30/09 a las 19:00 (t=219 h). \textbf{La campaña completa cierra el jueves 01/10/2026} (muestra condicional de F2 + desmontaje).

\section{Anclas temporales y definiciones}
\begin{itemize}\setlength{\itemsep}{1pt}
\item \textbf{PI (pre-inóculo):} muestra del mosto antes de inocular, lunes 15:30. Sin Oculyze. Cierra la deuda de $E_0$/YAN$_0$ y da la condición inicial del pool de CO$_2$ disuelto (en agosto el mosto arrancó con $\approx$430--470 mg/L, mientras el modelo hoy inicializa el pool en 0).
\item \textbf{t$_0$:} lunes 21/09/2026 \textbf{16:00}, fin del pulso de inoculación de la bomba peristáltica (programada 15:55$\rightarrow$16:00, $\approx$4--5 min). Todos los tiempos $t$ [h] se miden desde t$_0$.
\item El par PI/t$_0$ (15:30 / 16:10) queda a 40 min: única excepción a la separación mínima, por ser el par ancla pre/post inoculación (y duplicado de método para el Carbodoseur).
\item \textbf{Visitas y procesamiento:} cada visita demora \textbf{$\approx$1.5--2 h} (3 fermentadores $\times$ Y15, Oculyze, Carbodoseur, Brix, densidad, DO). Visitas a las \textbf{08:30} (listo $\sim$10:30) y \textbf{14:30} (listo $\sim$16:30, con margen antes de la limpieza de 17:00); viernes solo 08:30 (salida 12:30). Fin de semana sin personal.
\end{itemize}

\section{Perfiles de temperatura (programados en baños/controladores)}
\begin{table}[h]\centering
\begin{tabular}{L{3.4cm} L{1.8cm} L{2.6cm} L{1.8cm} L{3.1cm} L{7.6cm}}
\toprule
\rowcolor{primlight}\textbf{Reactor} & \textbf{Fase} & \textbf{Setpoint} & \textbf{t [h]} & \textbf{Fecha-hora} & \textbf{Nota} \\
\midrule
LAB019 (A, de LAB004) & 1 & 18 \textdegree C & 0 & lun 21/09 16:00 & Setpoint inicial; mosto cargado ya equilibrado \\LAB019 (A, de LAB004) & 2 & 21 \textdegree C & 99 & vie 25/09 19:00 & Automático, sin personal (réplica exacta del original) \\LAB019 (A, de LAB004) & 3 & 16 \textdegree C & 171 & lun 28/09 19:00 & Automático (réplica exacta del original) \\LAB020 (C, de LAB006) & 1 & 16 \textdegree C & 0 & lun 21/09 16:00 & Setpoint inicial \\LAB020 (C, de LAB006) & 2 & 18 \textdegree C & 123 & s\'ab 26/09 19:00 & Automático, fin de semana (réplica exacta) \\LAB020 (C, de LAB006) & 3 & 21 \textdegree C & 219 & mi\'e 30/09 19:00 & Automático (réplica exacta del original) \\LAB021 (A, de LAB009) & 1 & 18 \textdegree C & 0 & lun 21/09 16:00 & Setpoint inicial \\LAB021 (A, de LAB009) & 2 & 21 \textdegree C & 96 & vie 25/09 16:00 & Automático, con personal (réplica exacta) \\LAB021 (A, de LAB009) & 3 & 16 \textdegree C & 168 & lun 28/09 16:00 & Automático (en el original fue inferido, sin evento registrado) \\
\bottomrule
\end{tabular}
\caption{Los tiempos replican exactamente los ejecutados en LAB004--006 (Eventos\_operacion).}
\end{table}

\section{Nutrición (automatizada con bomba peristáltica)}
\begin{itemize}\setlength{\itemsep}{1pt}
\item \textbf{Pulso 1 (con la inoculación):} Springferm Xtrem 1.00 g + FDA 0.40 g por reactor, \emph{dentro del inóculo} (50\% del suplemento total; igual en A y C). Lunes 15:55$\rightarrow$16:00.
\item \textbf{Pulso 2 F3 (A, de LAB009):} igual dosis, t=48 h $\Rightarrow$ \textbf{mié 23/09 16:00}. El original no registró este pulso en el workbook: t=48 h según el calendario ICS (19/04), coherente con el cruce de densidad 1040 ($\sim$t=44 h). \emph{Contingencia:} si densidad(-07, mié 11:15) $\geq$1052 g/L, reprogramar a mié 19:00; si $\leq$1036 g/L, adelantar a mié 14:45.
\item \textbf{Pulso 2 F1 (A, de LAB004):} t=51 h $\Rightarrow$ \textbf{mié 23/09 19:00} (LAB004: densidad 1044.3 g/L en t=50 h). \emph{Contingencia:} si densidad(-08, mié 14:30) $\geq$1052 g/L $\rightarrow$ jue 08:30; si $\leq$1036 g/L $\rightarrow$ mié 18:00.
\item \textbf{Pulso 2 F2 (C, de LAB006):} t=75 h $\Rightarrow$ \textbf{jue 24/09 19:00} (LAB006: densidad 1040.0 g/L en t=74 h). \emph{Contingencia:} si densidad(-12, jue 14:30) $\geq$1048 g/L $\rightarrow$ vie 08:30; si $\leq$1032 g/L $\rightarrow$ jue 18:00.
\item Verificar en el log de la bomba (flanco de caída) la ejecución de cada pulso a la mañana siguiente.
\end{itemize}

\section{Plan de muestreo}
Mar--jue: \textbf{4 muestreos por día} (08:30 y 14:30 visitas completas con panel; 11:15 y 16:40 mini-muestreos solo instrumental, $\approx$15--20 min, sin Y15/Oculyze/Alcolyzer). Viernes: \textbf{2 muestreos} (08:30 completa + 11:00 mini). Separación 2.2--3.3 h intra-día, 15.8--18 h entre días. Análisis por visita completa: \textbf{Y15} (glucosa, fructosa, PAN, NH$_4\rightarrow$YAN, glicerol), \textbf{Alcolyzer} (etanol), \textbf{Oculyze} (queda en su software), \textbf{densidad + Brix + DO}. Semana siguiente: cola liviana (lun 28 $\times$2, mar 29 $\times$1, mié 30 $\times$2, jue 01 condicional).

\subsection*{Carbodoseur (CO$_2$ disuelto) — 11 lecturas fijas + 1 condicional}
\begin{itemize}\setlength{\itemsep}{1pt}
\item \textbf{Eventos:} PI, -01, -02 a -05 (mar: trayectoria del onset 16.5/19.25/22.5/24.7 h), -06 y -08 (mié: meseta y pre-pulso F3/F1), -10 (jue: post-pulso F1/F3), -14 (vie: post-pulso F2), -16 (lun 28: post-cambios T); -21 (jue 01/10) solo si se toma ($\approx$100 mL por lectura; la alícuota se devuelve al reactor).
\item \textbf{Por qué estos puntos:} IC del pool en PI/t0 (el modelo hoy parte en 0 y el mosto real arrancó con $\sim$450 mg/L); \textbf{cuatro puntos el martes} cubren la subida del pool hasta el onset esperado de A (24--30 h según loggers de abril; 22--24 h en LAB016--018), la ventana donde las variantes \emph{threshold} vs \emph{continuous release} más divergen; meseta (-06) para el nivel vs Csat; pre/post-pulso (-08/-10/-14) para la respuesta al pulso; -16 tras los cambios T para validar Csat(T).
\item \textbf{Protocolo:} tomar SIEMPRE la primera alícuota, directo del puerto y sin burbujear; anotar \emph{volumen remanente} (mL) y T de medición en la planilla de registro; el CO$_2$ disuelto se completa después con la tabla digitalizada. Leer DO en la misma visita. PI sale del contenedor de mosto (no consume volumen de reactores).
\item \textbf{Volumen (balance real):} la alícuota del Carbodoseur ($\approx$100 mL) \textbf{se devuelve al reactor}: no se pierde. Las únicas pérdidas no retornables por reactor son: Alcolyzer \textbf{50 mL por muestra} (F1/F3: mar--mié--jue $\approx$3 muestras = 150 mL; F2 = LAB020, la C fría: $\approx$5 muestras = 250 mL); Y15 1 mL y Oculyze 1 mL por visita completa; $\approx$2 mL por visita de mangueras/operación (también en minis). Pérdida total $\approx$0.2 L (F1/F3) y $\approx$0.3 L (F2): volumen final $\approx$1.7--1.8 L de los 2.0 L iniciales. Registrar el volumen acumulado solo como control.
\end{itemize}

{\footnotesize
\begin{longtable}{c L{2.2cm} L{1.6cm} c L{1.9cm} L{1.9cm} L{1.9cm} c c L{5.8cm}}
\toprule
\rowcolor{primlight}\textbf{N$^\circ$} & \textbf{Fecha-hora} & \textbf{Día} & \textbf{t [h]} & \textbf{LAB019 (F1)} & \textbf{LAB020 (F2)} & \textbf{LAB021 (F3)} & \textbf{Carbo} & \textbf{Sep.} & \textbf{Nota} \\
\midrule
\endfirsthead
\toprule
\rowcolor{primlight}\textbf{N$^\circ$} & \textbf{Fecha-hora} & \textbf{Día} & \textbf{t [h]} & \textbf{LAB019 (F1)} & \textbf{LAB020 (F2)} & \textbf{LAB021 (F3)} & \textbf{Carbo} & \textbf{Sep.} & \textbf{Nota} \\
\midrule
\endhead
1 & 21/09 15:30 & lun & -- & PI & PI & PI & SÍ & -- & Pre-inóculo: cierra la deuda de E0/YAN0 y da la IC del pool de CO2 disuelto ($\neq$0). No consume volumen de los reactores \\2 & 21/09 16:10 & lun & 0.17 & 01 & 01 & 01 & SÍ & 0.67* & Muestra t=0. Par ancla PI/t0 (40 min): excepción documentada a la separación mínima; YAN de t0 resuelve si el pulso 1 entra o no en la química; X0 por Oculyze \\3 & 22/09 08:30 & mar & 16.5 & 02 & 02 & 02 & SÍ & 16.33 & Pre-onset. Inicio de la trayectoria del pool del martes (16.5 h) \\4 & 22/09 11:15 & mar & 19.25 & 03 & 03 & 03 & SÍ & 2.75 & Mini-muestreo (\textasciitilde{}15-20 min): trayectoria del pool pre-onset (19.25 h) \\5 & 22/09 14:30 & mar & 22.5 & 04 & 04 & 04 & SÍ & 3.25 & Borde del onset esperado de A (referencias: 24-30 h según loggers de abril; 22-24 h en LAB016-018 a 20 \textdegree{}C) \\6 & 22/09 16:40 & mar & 24.67 & 05 & 05 & 05 & SÍ & 2.17 & En/frontera del onset de A (\textasciitilde{}24.7 h): el punto que discrimina acumulación previa (threshold) vs emisión temprana (continuous) \\7 & 23/09 08:30 & mié & 40.5 & 06 & 06 & 06 & SÍ & 15.83 & Meseta exponencial: nivel del pool vs Csat \\8 & 23/09 11:15 & mié & 43.25 & 07 & 07 & 07 & -- & 2.75 & Mini-muestreo de densidad (vigilar cruce de 1040 para los pulsos de la tarde) \\9 & 23/09 14:30 & mié & 46.5 & 08 & 08 & 08 & SÍ & 3.25 & Pre-pulso 2: F3 pulsa a las 16:00 (t=48) y F1 a las 19:00 (t=51). Contingencia densidad F1: esperada \textasciitilde{}1044 g/L (LAB004 t=50); si $\geq$1052 reprogramar a jue 08:30; si $\leq$1036 adelantar a mié 18:00 \\10 & 23/09 16:40 & mié & 48.67 & 09 & 09 & 09 & -- & 2.17 & 2.7 h post-pulso de F3 / pre-pulso de F1 (19:00): contraste entre fermentadores en la misma visita \\11 & 24/09 08:30 & jue & 64.5 & 10 & 10 & 10 & SÍ & 15.83 & Post-pulso: F1 a 13.5 h y F3 a 16.5 h del suyo; F2 pre-pulso (75 h = jue 19:00). Verificar logs de bomba del miércoles \\12 & 24/09 11:15 & jue & 67.25 & 11 & 11 & 11 & -- & 2.75 & Mini-muestreo de densidad (cruce de 1040 de F2) \\13 & 24/09 14:30 & jue & 70.5 & 12 & 12 & 12 & -- & 3.25 & Última muestra antes del pulso 2 de F2 (19:00). Contingencia densidad: esperada \textasciitilde{}1040 g/L (LAB006 t=74); si $\geq$1048 reprogramar a vie 08:30; si $\leq$1032 adelantar a jue 18:00 \\14 & 24/09 16:40 & jue & 72.67 & 13 & 13 & 13 & -- & 2.17 & 2.7 h pre-pulso de F2 (19:00) \\15 & 25/09 08:30 & vie & 88.5 & 14 & 14 & 14 & SÍ & 15.83 & Post-pulso de F2 (13.5 h) + A tardía. Verificar log del pulso de F2 y del cambio T de F3 (vie 16:00); F1 cambia a 21 \textdegree{}C hoy 19:00 \\16 & 25/09 11:00 & vie & 91 & 15 & 15 & 15 & -- & 2.5 & Última del día antes del cambio T de F1 (19:00); cierre de semana (salida 12:30) \\17 & 28/09 08:30 & lun & 160.5 & 16 & 16 & 16 & SÍ & 69.5 & Post-fin de semana: A seca esperada (t=120-146 h). Post-cambios T: F1 a 21 \textdegree{}C, F2 a 18 \textdegree{}C, F3 aún 21 \textdegree{}C (salta a 16 \textdegree{}C hoy 16:00). Descargar/respaldar datos \\18 & 28/09 14:30 & lun & 166.5 & 17 & 17 & 17 & -- & 6 & Confirmación de sequedad de A: F1/F3 secas si GF<2 g/L en -16/-17 (entonces vaciar el martes) \\19 & 29/09 08:30 & mar & 184.5 & 18 & 18 & 18 & -- & 18 & Confirmación A (si -17 no bastó) + cola de C. Si A confirmada: vaciar/desmontar F1 y F3 hoy \\20 & 30/09 08:30 & mié & 208.5 & 19 & 19 & 19 & -- & 24 & C en fase final (LAB006: GF=4.7 g/L a t=170 h) \\21 & 30/09 14:30 & mié & 214.5 & 20 & 20 & 20 & -- & 6 & C casi seca / pre-cambio T de F2 (21 \textdegree{}C hoy 19:00, t=219 h) \\22 & 01/10 08:30 & jue & 232.5 & --- & -21 & --- & cond & 18 & CONDICIONAL: solo F2 y solo si GF(-20) > 2 g/L; 13 h después del salto a 21 \textdegree{}C. Luego vaciado de F2 y cierre de campaña \\
\bottomrule
\end{longtable}
}
\footnotesize *Par ancla PI/t$_0$ (excepción documentada). Sep.\ = horas desde la visita anterior. Carbo = lectura de Carbodoseur (CO$_2$ disuelto); cond = solo si se toma -21.

\subsection*{Criterios de término}
\begin{itemize}\setlength{\itemsep}{1pt}
\item \textbf{A (F1/F3 = LAB019/021):} GF $<$ 2 g/L en dos muestras consecutivas (esperado -16/-17, lunes 28/09); el martes 29 por la mañana se vacían y desmontan F1/F3 si se confirma.
\item \textbf{C (F2 = LAB020):} GF $<$ 2 g/L en dos consecutivas (esperado -19/-20, miércoles 30/09); -21 (jueves 01/10 08:30) solo si GF(-20) $>$ 2 g/L; luego vaciado de F2 y cierre.
\item \textbf{Ventana ciega de fin de semana:} la sequedad de A (t=120--146 h) caerá entre el sábado 16:00 y el domingo 18:00 sin muestreo; la señal continua de CO$_2$ ---justo lo que esta réplica corrige--- cubre el intervalo.
\end{itemize}

\section{Cronograma día a día}
{\footnotesize
\renewcommand{\arraystretch}{1.15}
\begin{longtable}{L{1.7cm} L{1.2cm} L{1.1cm} c L{4.6cm} L{1.6cm} L{1.7cm} L{7.2cm}}
\toprule
\rowcolor{primlight}\textbf{Fecha} & \textbf{Día} & \textbf{Hora} & \textbf{t [h]} & \textbf{Evento} & \textbf{Reactores} & \textbf{Tipo} & \textbf{Detalle} \\
\midrule
\endfirsthead
\toprule
\rowcolor{primlight}\textbf{Fecha} & \textbf{Día} & \textbf{Hora} & \textbf{t [h]} & \textbf{Evento} & \textbf{Reactores} & \textbf{Tipo} & \textbf{Detalle} \\
\midrule
\endhead
21/09 & lun & 08:30 &  & Verificación del sistema CO2 sin fugas & F1,F2,F3 & Preparación & Prueba de presión/agua jabonosa en todas las conexiones y trampas; zero-check de sensores; comparar baseline F1 vs F2/F3 \\21/09 & lun & 10:00 &  & Calibración de bombas peristálticas & F1,F2,F3 & Preparación & Verificar caudal/volumen de las bombas de inóculo y nutrición; cargar programas de la semana \\21/09 & lun & 10:00 &  & Baños/controladores a setpoint inicial & F1,F2,F3 & Preparación & F1,F2: 18 \textdegree{}C; F3: 16 \textdegree{}C (para que el mosto llegue equilibrado a t0) \\21/09 & lun & 11:00 &  & Preparación de dosis, inóculo y material & F1,F2,F3 & Preparación & Pesar Springferm Xtrem (1.00 g/dosis) y FDA (0.40 g/dosis); rehidratar inóculo Zymaflore X5; tubos, etiquetas, hielo; Y15 y Alcolyzer operativos; imprimir planilla de registro; iniciar loggers CO2/T \\21/09 & lun & 14:30 &  & Carga y homogeneización del mosto & F1,F2,F3 & Preparación & Mosto natural Sauvignon blanc, 2 L por reactor; registrar volumen y temperatura \\21/09 & lun & 15:30 &  & MUESTRA PI (pre-inóculo) + Carbodoseur & F1,F2,F3 & Muestreo & Carbodoseur primero (\textasciitilde{}100 mL/reactor), luego panel. Y15 de PI puede procesarse mientras corre la inoculación \\21/09 & lun & 15:40 &  & Carga de inóculo + suplemento en línea de bomba & F1,F2,F3 & Preparación & El suplemento del pulso 1 va dentro del inóculo \\21/09 & lun & 15:55 &  & Inoculación automatizada (pulso 1) & F1,F2,F3 & Automático & Bomba peristáltica 15:55$\rightarrow$16:00 (\textasciitilde{}4-5 min, ref. LAB016-018: 232-246 s) \\21/09 & lun & 16:00 & 0 & t0 = fin de la inoculación & F1,F2,F3 & Automático & Ancla temporal de toda la campaña \\21/09 & lun & 16:10 & 0.17 & MUESTRA -01 (t=0) + Carbodoseur & F1,F2,F3 & Muestreo & Panel completo; lunes ajustado (PI y t0 en \textasciitilde{}2 h): si no alcanza, Y15 restante en la tarde-jueves \\21/09 & lun & 16:15 & 0.25 & Verificación inicial de trazas & F1,F2,F3 & Verificación & Confirmar que CO2/T están registrando y que la bomba logueó la inoculación; 17:00-17:30 limpieza diaria \\22/09 & mar & 08:30 & 16.5 & MUESTRA 02 02 (t=16.5 h) + Carbodoseur & F1,F2,F3 & Muestreo & Pre-onset. Inicio de la trayectoria del pool del martes (16.5 h) \\22/09 & mar & 11:15 & 19.25 & MINI-MUESTRA 03 (t=19.25 h) + Carbodoseur & F1,F2,F3 & Muestreo & Mini-muestreo (\textasciitilde{}15-20 min): trayectoria del pool pre-onset (19.25 h) \\22/09 & mar & 14:30 & 22.5 & MUESTRA 04 04 (t=22.5 h) + Carbodoseur & F1,F2,F3 & Muestreo & Borde del onset esperado de A (referencias: 24-30 h según loggers de abril; 22-24 h en LAB016-018 a 20 \textdegree{}C) \\22/09 & mar & 16:40 & 24.67 & MINI-MUESTRA 05 (t=24.67 h) + Carbodoseur & F1,F2,F3 & Muestreo & En/frontera del onset de A (\textasciitilde{}24.7 h): el punto que discrimina acumulación previa (threshold) vs emisión temprana (continuous) \\23/09 & mié & 08:30 & 40.5 & MUESTRA 06 06 (t=40.5 h) + Carbodoseur & F1,F2,F3 & Muestreo & Meseta exponencial: nivel del pool vs Csat \\23/09 & mié & 11:15 & 43.25 & MINI-MUESTRA 07 (t=43.25 h) & F1,F2,F3 & Muestreo & Mini-muestreo de densidad (vigilar cruce de 1040 para los pulsos de la tarde) \\23/09 & mié & 14:30 & 46.5 & MUESTRA 08 08 (t=46.5 h) + Carbodoseur & F1,F2,F3 & Muestreo & Pre-pulso 2: F3 pulsa a las 16:00 (t=48) y F1 a las 19:00 (t=51). Contingencia densidad F1: esperada \textasciitilde{}1044 g/L (LAB004 t=50); si $\geq$1052 reprogramar a jue 08:30; si $\leq$1036 adelantar a mié 18:00 \\23/09 & mié & 16:00 & 48 & Pulso nutricional 2 — F3 (réplica de LAB009) & F3 & Automático & Springferm Xtrem 1.00 g + FDA 0.40 g (t=48 h, vía ICS del original). Llegando de la visita de las 14:30 \\23/09 & mié & 16:40 & 48.67 & MINI-MUESTRA 09 (t=48.67 h) & F1,F2,F3 & Muestreo & 2.7 h post-pulso de F3 / pre-pulso de F1 (19:00): contraste entre fermentadores en la misma visita \\23/09 & mié & 19:00 & 51 & Pulso nutricional 2 — F1 (réplica de LAB004) & F1 & Automático & Springferm Xtrem 1.00 g + FDA 0.40 g (t=51 h; original: densidad 1040) \\24/09 & jue & 08:30 & 64.5 & MUESTRA 10 10 (t=64.5 h) + Carbodoseur & F1,F2,F3 & Muestreo & Post-pulso: F1 a 13.5 h y F3 a 16.5 h del suyo; F2 pre-pulso (75 h = jue 19:00). Verificar logs de bomba del miércoles \\24/09 & jue & 08:30 & 64.5 & Verificar log: pulsos 2 de F3 y F1 ejecutados & F3,F1 & Verificación & Revisar flancos de la bomba en el logger (análogo a nutricion\_activa de LAB016-018) \\24/09 & jue & 11:15 & 67.25 & MINI-MUESTRA 11 (t=67.25 h) & F1,F2,F3 & Muestreo & Mini-muestreo de densidad (cruce de 1040 de F2) \\24/09 & jue & 14:30 & 70.5 & MUESTRA 12 12 (t=70.5 h) & F1,F2,F3 & Muestreo & Última muestra antes del pulso 2 de F2 (19:00). Contingencia densidad: esperada \textasciitilde{}1040 g/L (LAB006 t=74); si $\geq$1048 reprogramar a vie 08:30; si $\leq$1032 adelantar a jue 18:00 \\24/09 & jue & 16:40 & 72.67 & MINI-MUESTRA 13 (t=72.67 h) & F1,F2,F3 & Muestreo & 2.7 h pre-pulso de F2 (19:00) \\24/09 & jue & 19:00 & 75 & Pulso nutricional 2 — F2 (réplica de LAB006) & F2 & Automático & Springferm Xtrem 1.00 g + FDA 0.40 g (t=75 h; original: densidad 1040) \\25/09 & vie & 08:30 & 88.5 & MUESTRA 14 14 (t=88.5 h) + Carbodoseur & F1,F2,F3 & Muestreo & Post-pulso de F2 (13.5 h) + A tardía. Verificar log del pulso de F2 y del cambio T de F3 (vie 16:00); F1 cambia a 21 \textdegree{}C hoy 19:00 \\25/09 & vie & 08:30 & 88.5 & Verificar log: pulso 2 de F2 ejecutado & F2 & Verificación & Ídem \\25/09 & vie & 11:00 & 91 & MINI-MUESTRA 15 (t=91 h) & F1,F2,F3 & Muestreo & Última del día antes del cambio T de F1 (19:00); cierre de semana (salida 12:30) \\25/09 & vie & 11:30 & 91.5 & Cierre de fin de semana & F1,F2,F3 & Verificación & Confirmar setpoint F3 21 \textdegree{}C ejecutado (16:00) y F1 21 \textdegree{}C programado (19:00); F2 cambia el sábado 19:00; capacidad de loggers hasta el lunes; 12:00-12:30 limpieza y salida \\25/09 & vie & 16:00 & 96 & Setpoint 21 \textdegree{}C — F3 & F3 & Automático & t=96 h (réplica exacta del original, ejecutado con personal presente) \\25/09 & vie & 19:00 & 99 & Setpoint 21 \textdegree{}C — F1 & F1 & Automático & t=99 h (réplica exacta del original, sin personal) \\\rowcolor{gris} 26/09 & sáb & 00:00 &  & Fin de semana sin personal & F1,F2,F3 & Sin personal & Sábado y domingo: logging continuo CO2/T; únicos eventos: setpoint F2 18 \textdegree{}C (sáb 19:00, automático). Ventana ciega de muestreo: la sequedad de A (t=120-146 h) ocurrirá aquí; el CO2 continuo sin fugas cubre el intervalo \\\rowcolor{gris} 26/09 & sáb & 19:00 & 123 & Setpoint 18 \textdegree{}C — F2 & F2 & Automático & t=123 h (réplica exacta del original, sin personal) \\28/09 & lun & 08:30 & 160.5 & MUESTRA 16 16 (t=160.5 h) + Carbodoseur & F1,F2,F3 & Muestreo & Post-fin de semana: A seca esperada (t=120-146 h). Post-cambios T: F1 a 21 \textdegree{}C, F2 a 18 \textdegree{}C, F3 aún 21 \textdegree{}C (salta a 16 \textdegree{}C hoy 16:00). Descargar/respaldar datos \\28/09 & lun & 08:30 & 160.5 & Verificar cambios T del finde + respaldo & F1,F2,F3 & Verificación & Confirmar setpoint 18 \textdegree{}C en F2 (sáb 19:00); F3 cambia a 16 \textdegree{}C hoy 16:00 y F1 a 16 \textdegree{}C hoy 19:00; descargar y respaldar datos \\28/09 & lun & 14:30 & 166.5 & MUESTRA 17 17 (t=166.5 h) & F1,F2,F3 & Muestreo & Confirmación de sequedad de A: F1/F3 secas si GF<2 g/L en -16/-17 (entonces vaciar el martes) \\28/09 & lun & 16:00 & 168 & Setpoint 16 \textdegree{}C — F3 & F3 & Automático & t=168 h (en el original fue inferido, sin evento registrado); lunes 16:00, con personal \\28/09 & lun & 19:00 & 171 & Setpoint 16 \textdegree{}C — F1 & F1 & Automático & t=171 h (réplica exacta del original) \\29/09 & mar & 08:30 & 184.5 & MUESTRA 18 18 (t=184.5 h) & F1,F2,F3 & Muestreo & Confirmación A (si -17 no bastó) + cola de C. Si A confirmada: vaciar/desmontar F1 y F3 hoy \\29/09 & mar & 08:30 & 184.5 & Si A confirmada seca (-16/-17): vaciar F1 y F3 & F1,F3 & Cierre & Criterio: GF<2 g/L en 2 muestras consecutivas (-16/-17); desmontaje de reactores liberados \\30/09 & mié & 08:30 & 208.5 & MUESTRA 19 19 (t=208.5 h) & F1,F2,F3 & Muestreo & C en fase final (LAB006: GF=4.7 g/L a t=170 h) \\30/09 & mié & 14:30 & 214.5 & MUESTRA 20 20 (t=214.5 h) & F1,F2,F3 & Muestreo & C casi seca / pre-cambio T de F2 (21 \textdegree{}C hoy 19:00, t=219 h) \\30/09 & mié & 19:00 & 219 & Setpoint 21 \textdegree{}C — F2 & F2 & Automático & t=219 h (réplica exacta del original) \\01/10 & jue & 08:30 & 232.5 & MUESTRA -21 (condicional, solo F2) + Carbodoseur & F2 & Muestreo & Solo si GF(-20) > 2 g/L; luego vaciado de F2 \\01/10 & jue & 09:30 & 233.5 & Cierre de campaña & F1,F2,F3 & Cierre & Vaciado/desmontaje de F2, descarga y respaldo final de datos, limpieza del laboratorio \\
\bottomrule
\end{longtable}
}

\section{Restricciones y contingencias}
\begin{itemize}\setlength{\itemsep}{1pt}
\item \textbf{Horarios y procesamiento:} lun--jue 8:00--17:30: visitas completas 08:30 y 14:30 ($\approx$1.5--2 h de proceso cada una) y mini-muestreos 11:15/16:40 ($\approx$15--20 min). Viernes 08:00--12:30 (completa 08:30 + mini 11:00). Fin de semana sin personal.
\item \textbf{Separación entre muestras:} mínimo 2 h (piso 1.5 h); el grid mar--jue usa 2.2--3.3 h (los minis son livianos, sin Y15). Única excepción: el par PI/t$_0$ del lunes.
\item \textbf{Mosto retrasado:} si no está listo a las 15:30, correr PI y t$_0$ (p.\,ej.\ mosto 16:00 $\rightarrow$ PI 16:00, t$_0$ 16:30) y desplazar \emph{todos} los eventos automáticos por el mismo desfase, manteniendo los tiempos $t$ [h].
\item \textbf{Fugas CO$_2$:} verificación con presión/jabón el lunes a.\,m.; revisar trazas a diario (caída a cero o asimetría F1 vs F2/F3 son señales de fuga).
\item \textbf{Carbodoseur:} primera alícuota de cada visita marcada, sin burbujear; anotar volumen remanente y T; CO$_2$ después con la tabla digitalizada; $\approx$100 mL por lectura \textbf{que se devuelve al reactor} (el «Vol.\ rem.» es solo control operativo). La pérdida real es pequeña: Alcolyzer 50 mL/muestra + Y15/Oculyze 1 mL c/u por visita completa + $\approx$2 mL/visita de mangueras $\Rightarrow$ $\approx$0.2 L (F1/F3) y $\approx$0.3 L (F2) en total; volumen final $\approx$1.7--1.8 L de 2.0 L.
\item \textbf{Señales de alarma en el lag (lecciones de abril):} trazas en 0.000 exactos = línea sospechosa de bloqueo (patrón F1); baseline creciente o saltos = sospecha de fuga (patrón F2/F3); lo esperable es $\sim$0.5--1 sccm plano por desorción.
\item \textbf{Trazabilidad:} nombrar canales LAB019/020/021 antes del lunes 16:00; registrar la hora exacta de cada flanco de bomba; descargar y respaldar datos el lunes 28 a.\,m.\ y al cierre.
\end{itemize}

\vspace{4pt}
{\sloppy
\noindent\begin{minipage}{\linewidth}
\textcolor{gray}{\footnotesize Fuentes: \texttt{mosto\_natural\_xthiol.xlsx} (hojas \texttt{Eventos\_operacion}, \texttt{Perfiles\_temperatura}, \texttt{LAB004..LAB006}); \texttt{Planilla\_Maestra\_ID\_2026\_14\_07\_26.xlsx}; \texttt{docs/natural\_must\_temporal\_audit\_LAB001\_018.md}; \texttt{docs/natural\_must\_model\_calibration\_and\_estimator\_readiness\_2026.md} (CO2sat\_scale en cota inferior; Carbodoseur LAB013--015). Registro de mediciones: \texttt{registro\_muestras\_LAB019\_021\_2026.xlsx}.}
\end{minipage}\par}
\end{document}
